% Faithful official CCC2022 article14 Sections6–7.
\section*{6 The plane test is a characterization}
In this section we show that the multiplicity code $\mathrm{MRM}(q,m,d,s)$ can be characterized by restrictions to planes. Let $\mathbb F_q$ be a field and let $m,s\leq d$ be positive integers. For $\boldsymbol a,\boldsymbol b,\boldsymbol c\in\mathbb F_q^m$\footnote{In this section, we omit the requirement that $\boldsymbol a,\boldsymbol b$ be independent. If some degenerate plane shows that a table is not a low degree polynomial, then some actual plane will too.} define: $\ell_{\boldsymbol a,\boldsymbol b,\boldsymbol c}:\mathbb F_q^2\to\mathbb F_q^m$ by:
\[\ell_{\boldsymbol a,\boldsymbol b,\boldsymbol c}(t,r)=\boldsymbol at+\boldsymbol br+\boldsymbol c.\]
From Lemma 18 and Definition 48 we see that:
\paragraph{Theorem 59 (Completeness).} Suppose $d<sq-1$. If a table $T\in\Sigma_{m,s}^{q^m}$ satisfies $T\in\mathrm{MRM}(q,m,d,s)$ then for all $\boldsymbol a,\boldsymbol b,\boldsymbol c\in\mathbb F_q^m$,
\[\phi_{(\boldsymbol a,\boldsymbol b)}\circ T\circ\ell_{\boldsymbol a,\boldsymbol b,\boldsymbol c}\in\mathrm{MRM}(q,2,d,s).\]
The main challenge is proving the converse:
\paragraph{Theorem 60 (Soundness).} Suppose $q$ is a power of the prime $p$, $q\geq s$ and $d<q(s-1/p)$. If a table $T\in\Sigma_{m,2}^{q^m}$ satisfies that for all $\boldsymbol a,\boldsymbol b,\boldsymbol c\in\mathbb F_q^m$, $\phi_{(\boldsymbol a,\boldsymbol b)}\circ T\circ\ell_{\boldsymbol a,\boldsymbol b,\boldsymbol c}\in\mathrm{MRM}(q,2,d,s)$ then $T\in\mathrm{MRM}(q,m,d,s)$.
We define the vector space of tables which pass the test:
\[V_{m,d,s}=\{T\in\Sigma_{m,s}^{q^m}\mid\forall\boldsymbol a,\boldsymbol b,\boldsymbol c\in\mathbb F_q^m\ \phi_{(\boldsymbol a,\boldsymbol b)}\circ T\circ\ell_{\boldsymbol a,\boldsymbol b,\boldsymbol c}\in\mathrm{MRM}(q,2,d,s)\}.\tag{12}\]
We denote by $C_{m,d,s}=V_{m,d,s}\setminus\mathrm{MRM}_{m,d,s}$, the set of tables which cheat the test. We would like to show that $C_{m,d,s}=\emptyset$. Assume towards a contradiction that there is a table $T\in C_{m,d,s}$. By Claim 41, $T$ can be realized (uniquely) as an element $P$ of the quotient space $\mathbb F_q[\boldsymbol X]/I_{m,s}$. We use the basis $B_{m,s}$ from Equation (5) to write $P$ in the form
\[P(\boldsymbol X)=\sum_{(\boldsymbol I,\boldsymbol J)\in M_{s,q}}\alpha_{\boldsymbol I,\boldsymbol J}\cdot g(\boldsymbol X)^{\boldsymbol I}\boldsymbol X^{\boldsymbol J},\]
where $\alpha_{\boldsymbol I,\boldsymbol J}\in\mathbb F_q$ and $(\boldsymbol I,\boldsymbol J)\in M_{s,q}$ iff $w(\boldsymbol I)<s$ and $J_k<q$ for every $1\leq k\leq m$. Since $T\notin\mathrm{MRM}_{m,d,s}$ we have $\deg(P)>d$. This means there must be some $\boldsymbol I$ and $\boldsymbol J$ such that $\alpha_{\boldsymbol I,\boldsymbol J}\neq0$ and
\[w(\boldsymbol I)q+w(\boldsymbol J)>d.\tag{13}\]
We may assume that every $\boldsymbol I,\boldsymbol J$ for which $\alpha_{\boldsymbol I,\boldsymbol J}\neq0$ satisfy Equation (13). This is since the test is linear, and any degree $\leq d$ terms have no effect on whether $P$ passes the test or not.
We use the following monomial order $\succ_w$ on $B_{m,s}$:
\begin{enumerate}
\item First order monomials according to $w(\boldsymbol I)+w(\boldsymbol J)$,
\item Then order monomials according to $w(\boldsymbol I)$,
\item Finally, order monomials according to the lexicographic order on $\boldsymbol I,\boldsymbol J$.
\end{enumerate}
We emphasize that $\succ_w$ is not a monomial order in the sense of Grobner bases, and we make no use of it in that sense.
Let $(\boldsymbol I_{\max},\boldsymbol J_{\max})$ be s.t. $g(\boldsymbol X)^{\boldsymbol I_{\max}}\cdot\boldsymbol X^{\boldsymbol J_{\max}}$ is a maximal monomial of $P$ according to $\succ_w$.
For $\boldsymbol a,\boldsymbol b,\boldsymbol c\in\mathbb F^m$ the restriction of $P$ to the plane defined by $\boldsymbol a,\boldsymbol b,\boldsymbol c$ is $P|_{P_{\boldsymbol a,\boldsymbol b,\boldsymbol c}}(t,r)=P(\boldsymbol at+\boldsymbol br+\boldsymbol c)$. Expressing $P|_{P_{\boldsymbol a,\boldsymbol b,\boldsymbol c}}(t,r)$ in the basis $B_{2,\infty}$:
\begin{align*}
P|_{P_{\boldsymbol a,\boldsymbol b,\boldsymbol c}}(t,r)&=\sum_{i\in\mathbb N,j\in\mathbb N,k,\ell<q}A_{i,j,k,\ell}(\boldsymbol a,\boldsymbol b,\boldsymbol c)\cdot g(t)^ig(r)^jt^kr^\ell,\\
P|_{P_{\boldsymbol a,\boldsymbol b,\boldsymbol c}}(t,r)\bmod I_{2,s}&=\sum_{i+j<s,k,\ell<q}A_{i,j,k,\ell}(\boldsymbol a,\boldsymbol b,\boldsymbol c)\cdot g(t)^ig(r)^jt^kr^\ell.
\end{align*}
We view $A_{i,j,k,\ell}(\boldsymbol a,\boldsymbol b,\boldsymbol c)$ as a polynomial in the variables $\boldsymbol a,\boldsymbol b,\boldsymbol c$.
\paragraph{Lemma 61.} For every partition $\boldsymbol J_{\max}=\boldsymbol J_{\max}^b+\boldsymbol J_{\max}^c$ such that:
\begin{itemize}
\item $q\cdot w(\boldsymbol I_{\max})+w(\boldsymbol J_{\max}^b)\leq qs-1$, and,
\item $\binom{\boldsymbol J_{\max}}{\boldsymbol J_{\max}^b}\neq0\bmod p$,
\end{itemize}
the monomial $\boldsymbol a^{\boldsymbol I_{\max}}\boldsymbol b^{\boldsymbol J_{\max}^b}\boldsymbol c^{\boldsymbol J_{\max}^c}$ appears with a non-zero coefficient at
\[A_{w(\boldsymbol I_{\max}),\lfloor w(\boldsymbol J_{\max}^b)/q\rfloor,0,w(\boldsymbol J_{\max}^b)\bmod q}(\boldsymbol a,\boldsymbol b,\boldsymbol c)\bmod I_{3m,1}.\]
Where the ideal $I_{3m,1}$ above is in the variables $a_1,\ldots,a_m,b_1,\ldots,b_m,c_1,\ldots,c_m$.
\paragraph{Proof.} We expand $P|_{P_{\boldsymbol a,\boldsymbol b,\boldsymbol c}}(t,r)$. First, $P(\boldsymbol X)=\sum_{(\boldsymbol I,\boldsymbol J)\in M_{s,q}}\alpha_{\boldsymbol I,\boldsymbol J}\cdot g(\boldsymbol X)^{\boldsymbol I}\boldsymbol X^{\boldsymbol J}$. Thus
\begin{align*}
P|_{P_{\boldsymbol a,\boldsymbol b,\boldsymbol c}}(t,r)
&=\sum_{(\boldsymbol I,\boldsymbol J)\in M_{s,q}}\alpha_{\boldsymbol I,\boldsymbol J}\cdot g(\boldsymbol at+\boldsymbol br+\boldsymbol c)^{\boldsymbol I}(\boldsymbol at+\boldsymbol br+\boldsymbol c)^{\boldsymbol J}\\
&=\sum_{(\boldsymbol I,\boldsymbol J)\in M_{s,q}}\alpha_{\boldsymbol I,\boldsymbol J}\prod_{k=1}^m(g(a_kt+b_kr+c_k))^{I_k}\cdot\prod_{i=k}^m(a_kt+b_kr+c_k)^{J_k}\\
&=\sum_{(\boldsymbol I,\boldsymbol J)\in M_{s,q}}\alpha_{\boldsymbol I,\boldsymbol J}\prod_{k=1}^m(g(t)a_k+g(r)b_k)^{I_k}\cdot\prod_{i=k}^m(a_kt+b_kr+c_k)^{J_k}\\
&=\sum_{(\boldsymbol I,\boldsymbol J)\in M_{s,q}}\alpha_{\boldsymbol I,\boldsymbol J}\sum_{\substack{\boldsymbol I_a+\boldsymbol I_b=\boldsymbol I\\\boldsymbol J_a+\boldsymbol J_b+\boldsymbol J_c=\boldsymbol J}}\binom{\boldsymbol I}{\boldsymbol I_a}\binom{\boldsymbol J}{\boldsymbol J_a,\boldsymbol J_b,\boldsymbol J_c}\\
&\qquad\boldsymbol a^{\boldsymbol I_a+\boldsymbol J_a}\boldsymbol b^{\boldsymbol I_b+\boldsymbol J_b}\boldsymbol c^{\boldsymbol J_c}\cdot g(t)^{w(\boldsymbol I_a)}g(r)^{w(\boldsymbol I_b)}t^{w(\boldsymbol J_a)}r^{w(\boldsymbol J_b)}.\tag{14}
\end{align*}
We expand $t^{w(\boldsymbol J_a)}$ and $r^{w(\boldsymbol J_b)}$ in the basis $B_{1,s}$ as in Example 43 to get:
\begin{align*}
P|_{P_{\boldsymbol a,\boldsymbol b,\boldsymbol c}}(t,r)
&=\sum_{(\boldsymbol I,\boldsymbol J)\in M_{s,q}}\alpha_{\boldsymbol I,\boldsymbol J}\cdot\sum_{\substack{\boldsymbol I_a+\boldsymbol I_b=\boldsymbol I\\\boldsymbol J_a+\boldsymbol J_b+\boldsymbol J_c=\boldsymbol J}}\binom{\boldsymbol I}{\boldsymbol I_a}\binom{\boldsymbol J}{\boldsymbol J_a\ \boldsymbol J_b\ \boldsymbol J_c}\\
&\quad\sum_{\substack{i_1,i_2\\i_1q+i_2\leq w(\boldsymbol J_a)\\i_2<q}}\sum_{\substack{j_1,j_2\\j_1q+j_2\leq w(\boldsymbol J_b)\\j_2<q}}\beta_{w(\boldsymbol J_a),i_1,i_2}\beta_{w(\boldsymbol J_b),j_1,j_2}\\
&\quad\boldsymbol a^{\boldsymbol I_a+\boldsymbol J_a}\boldsymbol b^{\boldsymbol I_b+\boldsymbol J_b}\boldsymbol c^{\boldsymbol J_c}\cdot g(t)^{w(\boldsymbol I_a)+i_1}t^{i_2}\cdot g(r)^{w(\boldsymbol I_b)+j_1}r^{j_2}.\tag{15}
\end{align*}
We now have a representation in the basis $B_{2,\infty}$.
Set:
\begin{align*}
\Delta_{rq}&=\left\lfloor\frac{w(\boldsymbol J_{\max}^b)}q\right\rfloor,\\
\Delta_t&=0,\\
\Delta_r&=w(\boldsymbol J_{\max}^b)\bmod q.\tag{16}
\end{align*}
We wish to see for which choice of values $\boldsymbol I,\boldsymbol J,\boldsymbol I_a,\boldsymbol I_b,\boldsymbol J_a,\boldsymbol J_b,\boldsymbol J_c,i_1,i_2,j_1,j_2$ in Equation (14), $\boldsymbol a^{\boldsymbol I_{\max}}\boldsymbol b^{\boldsymbol J_{\max}^b}\boldsymbol c^{\boldsymbol J_{\max}^c}$ appears as a coefficient of $A_{(w(\boldsymbol I_{\max}),\Delta_{rq},\Delta_t,\Delta_r)}\bmod I_{3m,1}$. We must have
\begin{align*}
\boldsymbol I_{\max}&=_{\mathbb F_q}\boldsymbol I_a+\boldsymbol J_a&&\text{By comparing the powers of }\boldsymbol a,\text{ remembering }\bmod I_{3m,1},\\
\boldsymbol J_{\max}^b&=_{\mathbb F_q}\boldsymbol I_b+\boldsymbol J_b&&\text{By comparing the powers of }\boldsymbol b,\text{ remembering }\bmod I_{3m,1},\\
\boldsymbol J_{\max}^c&=_{\mathbb F_q}\boldsymbol J_c&&\text{By comparing the powers of }\boldsymbol c,\text{ remembering }\bmod I_{3m,1},
\end{align*}
where $=_{\mathbb F_q}$ was defined in Definition 36.
By Claim 38 together with $s<q$ we have
\begin{align*}
w(\boldsymbol I_{\max})&\leq w(\boldsymbol I_a)+w(\boldsymbol J_a),\tag{17}\\
w(\boldsymbol J_{\max}^b)&\leq w(\boldsymbol I_b)+w(\boldsymbol J_b),\tag{18}\\
\boldsymbol J_{\max}^c&=\boldsymbol J_c.\tag{19}
\end{align*}
It follows that
\begin{align*}
w(\boldsymbol I_{\max})+w(\boldsymbol J_{\max})
&=w(\boldsymbol I_{\max})+w(\boldsymbol J_{\max}^b)+w(\boldsymbol J_{\max}^c)\\
&\leq w(\boldsymbol I_a)+w(\boldsymbol J_a)+w(\boldsymbol I_b)+w(\boldsymbol J_b)+w(\boldsymbol J_c)\\
&=w(\boldsymbol I)+w(\boldsymbol J).
\end{align*}
As $(\boldsymbol I_{\max},\boldsymbol J_{\max})$ is maximal it follows that
\[w(\boldsymbol I_{\max})+w(\boldsymbol J_{\max})=w(\boldsymbol I)+w(\boldsymbol J).\tag{20}\]
This, in turn, implies that both inequalities in Equations (17) and (18) are in fact equalities, i.e.,
\begin{align*}
w(\boldsymbol I_{\max})&=w(\boldsymbol I_a)+w(\boldsymbol J_a),\tag{21}\\
w(\boldsymbol J_{\max}^b)&=w(\boldsymbol I_b)+w(\boldsymbol J_b).\tag{22}
\end{align*}
We now look at
\[\deg_t\stackrel{\mathrm{def}}=\deg_t(g(t\boldsymbol a)^{\boldsymbol I_a}\cdot(t\boldsymbol a)^{\boldsymbol J_a}\cdot g(r\boldsymbol b)^{\boldsymbol I_b}\cdot(r\boldsymbol b)^{\boldsymbol J_b}\cdot\boldsymbol c^{\boldsymbol J_c}\bmod I_{2,s}).\]
On the one hand, we look for the monomial $\boldsymbol a^{\boldsymbol I_{\max}}\boldsymbol b^{\boldsymbol J_{\max}^b}\boldsymbol c^{\boldsymbol J_{\max}^c}$ in $A_{w(\boldsymbol I_{\max}),\Delta_{rq},\Delta_t,\Delta_r}\bmod I_{3m,1}$, and so we should have $\deg_t=q\cdot w(\boldsymbol I_{\max})+\Delta_t$. On the other hand, by Lemma 45,
\begin{align*}
\deg_t&=\deg_t(g(t\boldsymbol a)^{\boldsymbol I_a}(t\boldsymbol a)^{\boldsymbol J_a}g(r\boldsymbol b)^{\boldsymbol I_b}(r\boldsymbol b)^{\boldsymbol J_b}\boldsymbol c^{\boldsymbol J_c}\bmod I_{2,s})\\
&\leq\deg_t(g(t\boldsymbol a)^{\boldsymbol I_a}(t\boldsymbol a)^{\boldsymbol J_a})\\
&=q\cdot w(\boldsymbol I_a)+w(\boldsymbol J_a).
\end{align*}
Thus,
\[q\cdot w(\boldsymbol I_{\max})+\Delta_t\leq q\cdot w(\boldsymbol I_a)+w(\boldsymbol J_a).\]
Together with Equation (21) we see that $q\cdot w(\boldsymbol I_a)+q\cdot w(\boldsymbol J_a)+\Delta_t\leq q\cdot w(\boldsymbol I_a)+w(\boldsymbol J_a)$, and therefore $q\cdot w(\boldsymbol J_a)+\Delta_t\leq w(\boldsymbol J_a)$ which is possible iff $w(\boldsymbol J_a)=0$ (and $\Delta_t=0$, which is indeed true, see Equation (16)). Now, $w(\boldsymbol J_a)=0$ implies:
\[w(\boldsymbol I_a)=w(\boldsymbol I_{\max}),\text{ and,}\qquad\boldsymbol J_a=\emptyset.\]
We saw in Equation (20) that $w(\boldsymbol I)+w(\boldsymbol J)=w(\boldsymbol I_{\max})+w(\boldsymbol J_{\max})$. If $\boldsymbol I_b\neq\emptyset$ then
\[w(\boldsymbol I)=w(\boldsymbol I_a)+w(\boldsymbol I_b)>w(\boldsymbol I_a)=w(\boldsymbol I_a)+w(\boldsymbol J_a)=w(\boldsymbol I_{\max}).\]
Thus, $(\boldsymbol I,\boldsymbol J)\succ_w(\boldsymbol I_{\max},\boldsymbol J_{\max})$ in contradiction to the maximality of $(\boldsymbol I_{\max},\boldsymbol J_{\max})$. We conclude that $\boldsymbol I_b=\emptyset$.
As $(\boldsymbol I_{\max},\boldsymbol J_{\max})\in M_{s,q}$ we have $(\boldsymbol J_{\max})_k\leq q-1$ for every $k\in[m]$. Thus, $\boldsymbol J_{\max}^b=_{\mathbb F_q}\boldsymbol J_b$, $w(\boldsymbol J_{\max}^b)=w(\boldsymbol J_b)$ and $\boldsymbol J_{\max}^b$ is already reduced. Together this implies that $\boldsymbol J_{\max}^b=\boldsymbol J_b$.
Finally, we use the hypothesis that $q\geq s$. We have, $(\boldsymbol I_{\max})_k<s\leq q$ for all $k\in[m]$. Thus, $\boldsymbol I_{\max}=_{\mathbb F_q}\boldsymbol I_a$, $w(\boldsymbol I_{\max})=w(\boldsymbol I_a)$ and $\boldsymbol I_{\max}$ is already reduced. Together this implies that that $\boldsymbol I_{\max}=\boldsymbol I_a$.
Thus,
\[\boldsymbol I_{\max}=\boldsymbol I_a=\boldsymbol I_a+\boldsymbol I_b=\boldsymbol I,\qquad\boldsymbol J_{\max}^b=\boldsymbol I_b+\boldsymbol J_b=\boldsymbol J_b.\]
Altogether,the only term that may possibly contribute $\boldsymbol a^{\boldsymbol I_{\max}}\boldsymbol b^{\boldsymbol J_{\max}^b}\boldsymbol c^{\boldsymbol J_{\max}^c}$ to $A_{w(\boldsymbol I_{\max}),\Delta_{rq},\Delta_t,\Delta_r}\bmod I_{3m,1}$ is
\[(\boldsymbol I_a,\boldsymbol I_b,\boldsymbol J_a,\boldsymbol J_b,\boldsymbol J_c)=(\boldsymbol I_{\max},\emptyset,\emptyset,\boldsymbol J_{\max}^b,\boldsymbol J_{\max}^c).\]
Also, the tuple $(\boldsymbol I_a,\boldsymbol I_b,\boldsymbol J_a,\boldsymbol J_b,\boldsymbol J_c)=(\boldsymbol I_{\max},\emptyset,\emptyset,\boldsymbol J_{\max}^b,\boldsymbol J_{\max}^c)$ contributes
\begin{multline*}
\alpha_{\boldsymbol I_{\max},\boldsymbol J_{\max}}\cdot\binom{\boldsymbol J_{\max}}{\boldsymbol J_{\max}^b}\sum_{\substack{j_1,j_2\\j_1q+j_2\leq w(\boldsymbol J_{\max}^b)\\j_2<q}}\beta_{w(\boldsymbol J_{\max}^b),j_1,j_2}\\
\boldsymbol a^{\boldsymbol I_{\max}}\boldsymbol b^{\boldsymbol J_{\max}^b}\boldsymbol c^{\boldsymbol J_{\max}^c}\cdot g(t)^{w(\boldsymbol I_{\max})}\cdot g(r)^{j_1}r^{j_2}
\end{multline*}
to the term in Equation (14).
Notice that $w(\boldsymbol I_{\max})+j_1<s$, for otherwise $q\cdot w(\boldsymbol I_{\max})+w(\boldsymbol J_{\max}^b)\geq q(w(\boldsymbol I_{\max})+j_1)\geq qs$ in contradiction to the hypothesis. Thus the term is already $I_{2,s}$ reduced. The contribution to $A_{w(\boldsymbol I_{\max}),\Delta_{rq},0,\Delta_r}\bmod I_{3m,1}$ occurs for $(j_1,j_2)$ such that $j_1=\Delta_{rq}$ and $j_2=\Delta_r$. Thus, in the sum in Equation (14) there is exactly one possible way to contribute $\boldsymbol a^{\boldsymbol I_{\max}}\boldsymbol b^{\boldsymbol J_{\max}^b}\boldsymbol c^{\boldsymbol J_{\max}^c}$ to $A_{w(\boldsymbol I_{\max}),\Delta_{rq},0,\Delta_r}\bmod I_{3m,1}$, and this is when
\begin{align*}
(\boldsymbol I_a,\boldsymbol I_b,\boldsymbol J_a,\boldsymbol J_b,\boldsymbol J_c)&=(\boldsymbol I_{\max},\emptyset,\emptyset,\boldsymbol J_{\max}^b,\boldsymbol J_{\max}^c),\text{ and}\\
(i_1,i_2,j_1,j_2)&=(0,0,\Delta_{rq},\Delta_r).
\end{align*}
The coefficient of this term is:
\[\alpha_{\boldsymbol I_{\max},\boldsymbol J_{\max}}\cdot\binom{\boldsymbol J_{\max}}{\boldsymbol J_{\max}^b}\cdot\beta_{w(\boldsymbol J_{\max}^b),\Delta_{rq},\Delta_r}.\]
$\alpha_{\boldsymbol I_{\max},\boldsymbol J_{\max}}\neq0$. By assumption $\binom{\boldsymbol J_{\max}}{\boldsymbol J_{\max}^b}$ is non-zero. $\beta_{\ell,\lfloor\ell/q\rfloor,\ell\bmod q}=1$ (see Example 43) and taking $\ell=w(\boldsymbol J_{\max}^b)$ shows the coefficient is non-zero. As there is a unique term contributing the monomial with a non-zero coefficient, the monomial cannot cancel in $A_{w(\boldsymbol I_{\max}),\Delta_{rq},\Delta_t,\Delta_r}\bmod I_{3m,1}$ and $A_{w(\boldsymbol I_{\max}),\Delta_{rq},\Delta_t,\Delta_r}\bmod I_{3m,1}$ is non-zero.\hfill$\triangleleft$
\paragraph{Lemma 62.} There is a partition $\boldsymbol J_{\max}=\boldsymbol J_{\max}^b+\boldsymbol J_{\max}^c$, such that
\begin{itemize}
\item $d<q\cdot w(\boldsymbol I_{\max})+w(\boldsymbol J_{\max}^b)\leq qs-1$, and,
\item $\binom{\boldsymbol J_{\max}}{\boldsymbol J_{\max}^b}\bmod p\neq0$.
\end{itemize}
\paragraph{Proof.} Suppose $q=p^w$ where $p$ is prime (if $q$ is prime then $p=q$ and $w=1$). We choose $\boldsymbol J_{\max}^b$ as follows. We go over $k=1,\ldots,m$ and find the first $k_0\geq0$ such that $q\cdot w(\boldsymbol I_{\max})+\sum_{k=0}^{k_0}(\boldsymbol J_{\max})_k>d$. There must be some $k_0\leq m$ like that since $qw(\boldsymbol I_{\max})+w(\boldsymbol J_{\max})>d$. We set:
\begin{itemize}
\item $(\boldsymbol J_{\max}^b)_k=(\boldsymbol J_{\max})_k$, for $k=1,\ldots,k_0-1$, and,
\item $(\boldsymbol J_{\max}^b)_k=0$, for $k=k_0+1,\ldots,m$.
\end{itemize}
Set $v=(\boldsymbol J_{\max})_{k_0}$. There is a first value $0<v'\leq v$ such that
\[qw(\boldsymbol I_{\max})+\sum_{k=1}^{k_0-1}(\boldsymbol J_{\max})_k+v'>d.\]
We express $v=(\boldsymbol J_{\max})_{k_0}$ in base $p$: $v=\sum_{\ell=0}^{w-1}v_\ell\cdot p^\ell$. We let $v''$ be the first integer such that:
\begin{itemize}
\item $v''\geq v'$, and,
\item If we express $v''$ in base $p$ as $v''=\sum_{\ell=0}^{w-1}v''_\ell\cdot p^\ell$ then $v''_\ell\leq v_\ell$ for every $\ell$.
\end{itemize}
\paragraph{Claim 63.} $v'\leq v''\leq\min\{v,v'+p^{w-1}-1\}$.
\paragraph{Proof.} Notice that $v$ respects the conditions that we need, and so if $v\leq v'+p^{w-1}-1$ the claim holds. Otherwise, $v\geq v'+p^{w-1}-1$. Then, $v''\leq v'+p^{w-1}-1$ because we can increase $v'$ by setting all the lower bits in the $p$-representation to $0$, while increasing the most significant bit (that is multiplied by $p^{w-1}$) by $1$.\hfill$\triangleleft$
Thus, by Lucas theorem (see Equation (4))
\[\binom v{v''}\bmod p=\prod_{\ell=0}^{w-1}\binom{v_\ell}{v''_\ell}\neq0.\]
Having that, we let $(\boldsymbol J_{\max}^b)_{k_0}=v''$. We have, $\boldsymbol J_{\max}^b\leq\boldsymbol J_{\max}$, $qw(\boldsymbol I_{\max})+w(\boldsymbol J_{\max}^b)>d$ and $\binom{\boldsymbol J_{\max}}{\boldsymbol J_{\max}^b}\bmod p\neq0$. Also,
\[qw(\boldsymbol I_{\max})+w(\boldsymbol J_{\max}^b)\leq d+p^{w-1}\leq qs-1,\]
because $d\leq qs-p^{w-1}-1=q(s-1/p)-1$, completing the proof of the lemma.\hfill$\triangleleft$
We are now ready to prove Theorem 60.
\paragraph{Proof.} Fix a partition $\boldsymbol J_{\max}=\boldsymbol J_{\max}^b+\boldsymbol J_{\max}^c$ as in Lemma 62. Let
\begin{align*}
\deg_r&=\deg_r(g(t)^{w(\boldsymbol I_{\max})}r^{w(\boldsymbol J_{\max}^b)}\bmod I_{2,s})\\
&=\deg_r(g(t)^{w(\boldsymbol I_{\max})}r^{w(\boldsymbol J_{\max}^b)})=w(\boldsymbol J_{\max}^b).
\end{align*}
Define $\Delta_{rq}=\lfloor w(\boldsymbol J_{\max}^b)/q\rfloor$, $\Delta_t=0$ and $\Delta_r=w(\boldsymbol J_{\max}^b)\bmod q$. By Lemma 61 we know that
\[A_{w(\boldsymbol I_{\max}),\Delta_{rq},\Delta_t,\Delta_r}(\boldsymbol a,\boldsymbol b,\boldsymbol c)\bmod I_{3m,1}\neq0.\]
Thus, there exist $\boldsymbol a_0,\boldsymbol b_0,\boldsymbol c_0$ such that
\[A_{w(\boldsymbol I_{\max}),\Delta_{rq},\Delta_t,\Delta_r}(\boldsymbol a_0,\boldsymbol b_0,\boldsymbol c_0)\neq0.\]
We look at the test $\boldsymbol a_0,\boldsymbol b_0,\boldsymbol c_0$. We have
\[P|_{P_{\boldsymbol a,\boldsymbol b,\boldsymbol c}}(r,t)\bmod I_{2,s}=\sum_{i+j<s,k,\ell<q}A_{i,j,k,\ell}(\boldsymbol a,\boldsymbol b,\boldsymbol c)g(t)^ig(r)^jt^kr^\ell.\]
As $A_{w(\boldsymbol I_{\max}),\Delta_{rq},\Delta_t,\Delta_r}(\boldsymbol a_0,\boldsymbol b_0,\boldsymbol c_0)\neq0$ we see that
\begin{align*}
\deg(P|_{P_{\boldsymbol a_0,\boldsymbol b_0,\boldsymbol c_0}}\bmod I_{2,s})
&\geq q\cdot w(\boldsymbol I_{\max})+q\cdot\Delta_{rq}+\Delta_t+\Delta_r\\
&=q\cdot w(\boldsymbol I_{\max})+w(\boldsymbol J_{\max}^b)>d.
\end{align*}
Thus, by Lemma 47, the test $(\boldsymbol a_0,\boldsymbol b_0,\boldsymbol c_0)$ rejects.\hfill$\triangleleft$
\section*{7 Planes give a local MRM tester}
We restate Theorem 9:
\paragraph{Theorem 64.} Suppose $q$ is a prime power, $s\leq q$ and $d<q(s-1/p)$. Let $T:\mathbb F_q^m\to\Sigma_{m,s}$ be a table and let $\delta=\delta(T,\mathrm{MRM}(q,m,d,s))$. Then
\[\mathrm{REJ}_{2,d}^{\mathrm{MRM}}(T)\geq\min\{\alpha\delta,c\}\]
with $\alpha=\Omega(q^{-6s+5})$ and $c=\Omega(q^{-8s+4})$.
\paragraph{Proof.} We remind the reader that
\[H_k=\{\boldsymbol h=(\boldsymbol h_0,\boldsymbol h_1,\ldots,\boldsymbol h_k)\mid\boldsymbol h_0,\ldots,\boldsymbol h_k\in\mathbb F_q^m,\dim(\operatorname{span}\{\boldsymbol h_1,\ldots,\boldsymbol h_k\})=k\}.\]
Let us say that $\boldsymbol v\in H_k$ is bad, if the $k$-dimensional test with $\boldsymbol v$ rejects, i.e., $(\phi_{\boldsymbol v}\circ T)|_{A_{\boldsymbol v}}\notin\mathrm{MRM}(q,k,d,s)$. We let $t=\lceil(d+1)/(q-q/p)\rceil$ be the $\mathrm{RM}(q,m,d,s)$ testing dimension. Selecting a uniform $\boldsymbol u\in H_2$ is the same as selecting $\boldsymbol h\sim H_t$ and then a uniform $\boldsymbol u\in H_2$ such that $A_{\boldsymbol u}\subset A_{\boldsymbol h}$. Thus,
\begin{align*}
\mathrm{REJ}_{2,d}^{\mathrm{MRM}}(T)&=\Pr_{\boldsymbol u\sim H_2}(\boldsymbol u\text{ is bad})\\
&\geq\Pr_{\boldsymbol h\in H_t}(\boldsymbol h\text{ is bad})\cdot\Pr_{\boldsymbol u\in H_2:A_{\boldsymbol u}\subset A_{\boldsymbol h}}(\boldsymbol u\text{ is bad}\mid\boldsymbol h\text{ is bad}).
\end{align*}
\begin{itemize}
\item By Corollary 52, the probability of picking a bad $\boldsymbol h\in H_t$ is:
\[\mathrm{REJ}_{t,d}^{\mathrm{MRM}}(T)\geq\min\left\{\frac\delta{3q},\frac1{2(t+1)q^{t+1}}\right\}.\]
\item By Theorem 60 we know that for any bad $\boldsymbol h\in H_t$ there is at least one $\boldsymbol u\in H_2$ such that $A_{\boldsymbol u}$ is contained $A_{\boldsymbol h}$ and $\boldsymbol u$ is bad. Furthermore, if $A_{\boldsymbol u}=A_{\boldsymbol u'}$ then $\boldsymbol u$ is bad iff $\boldsymbol u'$ is bad. We look at $A_{\boldsymbol u}$. There are $(q^2-1)(q^2-q)q^2$ different $\boldsymbol u'\in H_2$ such that $A_{\boldsymbol u'}=A_{\boldsymbol u}$ (because there are $q^2-1$ choices for the first basis element, $q^2-q$ choices for the second basis element and $|A_{\boldsymbol u}|=q^2$ choices for the offset).
\end{itemize}
Altogether,
\[\mathrm{REJ}_{2,d}^{\mathrm{MRM}}(T)\geq\frac{\Omega(q^6)}{q^{3t}}\cdot\min\left\{\frac\delta{3q},\frac1{2(t+1)q^{t+1}}\right\}=\Omega\left(\min\left\{q^{-3t+5}\delta,\frac{q^{-4t+5}}t\right\}\right).\]
We notice that $t\leq\lceil q(s-1/p)/(q-q/p)\rceil\leq\lceil s/(1-1/p)\rceil\leq2s$. Thus $\alpha=\Omega(q^{-6s+5})$ and $c=\Omega(q^{-8s+4})$ as promised.\hfill$\triangleleft$
